% Fuente única de clase: Recursos/2025-11-04.pdf, páginas 1--9.
En la clase 20 retomamos los módulos libres para estudiar independencia lineal, bases y rango. Distinguimos después libertad y ausencia de torsión, y probamos un lema de separación que se utilizará para los submódulos sobre dominios de ideales principales.

%%%%%%%%%%%%%%%%%%%%%%%%%%%%%%%%%%%%%%%%%%%%%%%%%%%%%%%%%%%%%%%%%%%%%%%%%%%%%%%
\section{Módulos libres, independencia y bases}
%%%%%%%%%%%%%%%%%%%%%%%%%%%%%%%%%%%%%%%%%%%%%%%%%%%%%%%%%%%%%%%%%%%%%%%%%%%%%%%
%%%%%%%%%%%%%%%%%%%%%%%%%%%%%%%%%%%%%%%%%%%%%%%%%%%%%%%%%%%%%%%%%%%%%%%%%%%%%%%
\subsection{Familias y combinaciones lineales finitas}
%%%%%%%%%%%%%%%%%%%%%%%%%%%%%%%%%%%%%%%%%%%%%%%%%%%%%%%%%%%%%%%%%%%%%%%%%%%%%%%
% P01--P02.
En esta clase $R$ es conmutativo y los módulos son unitarios. Recordemos que un $R$-módulo es un grupo abeliano $M$ con una acción que satisface
\[
(r+s)m=rm+sm,\quad r(m+n)=rm+rn,\quad
(rs)m=r(sm),\quad 1m=m.
\]
Para un conjunto $I$, el módulo libre de la clase 14 es
\[
F^R(I)=R^{\oplus I}
=\set*{a:I\to R\given \supp(a)\text{ es finito}},
\qquad \supp(a)=\set{i\in I\given a(i)\ne0}.
\]
La función canónica $j:I\to R^{\oplus I}$ envía $i$ a $e_i$, donde $e_i(h)=1$ si $h=i$ y $e_i(h)=0$ en otro caso. Toda familia $u:I\to M$, con $m_i=u(i)$, induce un único homomorfismo
\[
\begin{tikzcd}[column sep=large,row sep=large]
R^{\oplus I}\arrow[r,"\varphi",destacado]&M\\
I\arrow[u,"j"]\arrow[ur,"u"']&
\end{tikzcd}
\qquad
\varphi\left(\sum_{i\in I}a_ie_i\right)=\sum_{i\in I}a_im_i.
\]
Ambas sumas tienen soporte finito, incluso cuando $I$ es infinito.

\begin{obs}
La función canónica $j$ es inyectiva si $R\ne0$. Para $R=0$, todo módulo unitario es cero y una familia indexada puede repetir elementos. Al identificar una base con un subconjunto supondremos $R\ne0$; dominios y campos ya cumplen esta condición.
\end{obs}

\begin{defi}
La familia $(m_i)_{i\in I}$ es \emph{linealmente independiente} si $\varphi$ es inyectiva, y \emph{dependiente} en caso contrario. Es \emph{generadora} si $\varphi$ es sobreyectiva y es una \emph{base} si $\varphi$ es un isomorfismo.
\end{defi}
La inyectividad equivale al criterio usual: para todo subconjunto finito $J\subseteq I$,
\[
\sum_{i\in J}a_im_i=0\quad\Longrightarrow\quad a_i=0\ (i\in J).
\]
En efecto, las relaciones de este tipo son precisamente los elementos de $\ker\varphi$, escritos en las coordenadas $e_i$. La generación significa que cada elemento de $M$ es una combinación lineal finita de la familia. Así,
\[
\ang{m_i\given i\in I}
=\set*{\sum_{i\in J}a_im_i\given J\subseteq I\text{ finito},\ a_i\in R}.
\]

\begin{ejem}
El subconjunto $\{2\}\subseteq\mathbb Z$ es linealmente independiente como familia de un $\mathbb Z$-módulo. Es maximal con esa propiedad: si añadimos $n\ne2$, la relación
\[
(-n)\,2+2\,n=0
\]
tiene coeficientes no todos nulos. Para $n=5$, obtenemos $(-5)2+2(5)=0$. Sin embargo, $\ang{2}=2\mathbb Z\ne\mathbb Z$, por lo que esa familia maximal no es una base.
\end{ejem}

\begin{lema}\label{lema:C20-libre-base}
Un módulo es libre si y sólo si admite una base. La familia $B\to M$ es una base exactamente cuando su homomorfismo inducido $R^{\oplus B}\to M$ es un isomorfismo.
\end{lema}
\begin{proof}
Ser base equivale a que ese homomorfismo sea inyectivo y sobreyectivo. Recíprocamente, un isomorfismo $R^{\oplus B}\to M$ transporta la familia canónica a una base de $M$. Cuando $R\ne0$, sus elementos son distintos y podemos identificar $B$ con su imagen.
\end{proof}

%%%%%%%%%%%%%%%%%%%%%%%%%%%%%%%%%%%%%%%%%%%%%%%%%%%%%%%%%%%%%%%%%%%%%%%%%%%%%%%
\subsection{Bases de espacios vectoriales}
%%%%%%%%%%%%%%%%%%%%%%%%%%%%%%%%%%%%%%%%%%%%%%%%%%%%%%%%%%%%%%%%%%%%%%%%%%%%%%%
% P02--P03.
Un módulo sobre un campo $k$ es un \emph{espacio vectorial}; sus elementos son vectores y los elementos de $k$ son escalares. La posibilidad de dividir por un escalar no nulo cambia el comportamiento del ejemplo anterior.

\begin{lema}\label{lema:C20-maximal-base}
Todo subconjunto linealmente independiente maximal de un espacio vectorial $V$ es una base de $V$.
\end{lema}
\begin{proof}
Sea $B$ maximal y $v\notin B$. La familia $B\cup\{v\}$ es dependiente, así que hay una relación finita no trivial
\[
cv+\sum_{b\in J}c_bb=0,\qquad J\subseteq B.
\]
El coeficiente $c$ es no nulo: de lo contrario tendríamos una relación no trivial entre elementos de $B$. Por tanto,
\[
v=-c^{-1}\sum_{b\in J}c_bb\in\ang{B}.
\]
Los elementos de $B$ también pertenecen a su generado; de aquí $V=\ang{B}$.
\end{proof}

\begin{prop}\label{prop:C20-extension-base}
Todo subconjunto linealmente independiente $S$ de un espacio vectorial $V$ se extiende a una base de $V$. En particular, todo espacio vectorial es un módulo libre.
\end{prop}
\begin{proof}
Consideremos el conjunto ordenado por inclusión
\[
\mathcal L=\set{A\subseteq V\given S\subseteq A,\ A\text{ linealmente independiente}}.
\]
Es no vacío, pues contiene a $S$. La unión de una cadena no vacía es independiente: una relación utiliza sólo finitos vectores, que pertenecen simultáneamente a un miembro de la cadena. La cadena vacía tiene a $S$ como cota. Por Zorn existe un elemento maximal $B$ de $\mathcal L$; cualquier extensión independiente de $B$ seguiría conteniendo a $S$, de modo que $B$ es maximal entre todos los subconjuntos independientes de $V$. Aplicamos el lema anterior y el \cref{lema:C20-libre-base}.
\end{proof}

%%%%%%%%%%%%%%%%%%%%%%%%%%%%%%%%%%%%%%%%%%%%%%%%%%%%%%%%%%%%%%%%%%%%%%%%%%%%%%%
\section{Cardinalidad de las bases y rango}
%%%%%%%%%%%%%%%%%%%%%%%%%%%%%%%%%%%%%%%%%%%%%%%%%%%%%%%%%%%%%%%%%%%%%%%%%%%%%%%
%%%%%%%%%%%%%%%%%%%%%%%%%%%%%%%%%%%%%%%%%%%%%%%%%%%%%%%%%%%%%%%%%%%%%%%%%%%%%%%
\subsection{Comparación de familias independientes}
%%%%%%%%%%%%%%%%%%%%%%%%%%%%%%%%%%%%%%%%%%%%%%%%%%%%%%%%%%%%%%%%%%%%%%%%%%%%%%%
% P03--P05.
\begin{prop}\label{prop:C20-cardinal}
Sea $R$ un dominio y $M$ un $R$-módulo libre. Si $B$ es un subconjunto linealmente independiente maximal de $M$ y $S$ es linealmente independiente, entonces $|S|\leq|B|$. En consecuencia, dos subconjuntos independientes maximales de $M$ tienen la misma cardinalidad.
\end{prop}
\begin{proof}
\emph{Paso al campo de fracciones.} Elijamos una base $E$ de $M$ e identifiquemos $M$ con $R^{\oplus E}\subseteq K^{\oplus E}=V$, donde $K=K(R)$. Una familia de elementos de $M$ es independiente sobre $R$ si y sólo si lo es sobre $K$: una relación sobre $R$ es también una relación sobre $K$; una relación finita sobre $K$ se convierte en una sobre $R$ multiplicando por el producto no nulo de sus denominadores. Si los coeficientes resultantes son cero, también lo eran los originales.

Además, $B$ genera $V$ sobre $K$. Para $m\in M\setminus B$, la maximalidad da
\[
cm+\sum_{b\in J}c_bb=0,\qquad c\ne0,
\]
y permite despejar $m$ en el $K$-generado de $B$. Como $M$ genera $V$ sobre $K$, resulta que $B$ es una base de $V$. Basta entonces comparar una familia independiente $S$ con una base $B$ de un espacio vectorial.

\emph{Intercambio finito.} Supongamos reemplazados algunos elementos distintos de $B$ por los vectores ya considerados de $S$, conservando una base $B'$. Para el siguiente vector $v\in S$ escribimos
\[
v=\sum_{s\text{ ya considerado}}c_ss
  +\sum_{b\in B\text{ no reemplazado}}c_bb.
\]
Algún $c_b$ es no nulo, pues $S$ es independiente. Elegimos uno y reemplazamos $b$ por $v$. Podemos despejar
\[
b=c_b^{-1}\left(v-\sum_s c_ss-\sum_{b'\ne b}c_{b'}b'\right).
\]
Así, la nueva familia genera; también es independiente, porque al sustituir $v$ en una relación, el coeficiente de $b$ obliga a anular el de $v$, y luego todos los demás. Este proceso asigna a cada vector considerado un elemento distinto de $B$. Si $B$ tiene $n$ elementos, no pueden existir $n+1$ vectores independientes: después de $n$ reemplazos, éstos ya generan $V$. Concluimos $|S|\leq n$.

\emph{Base infinita.} Para cada $s\in S$, sea $E_s\subseteq B$ el soporte finito de su expresión en la base. Para un conjunto finito fijo $E_0\subseteq B$, los vectores con $E_s=E_0$ son independientes en el espacio generado por $E_0$; el caso finito muestra que hay a lo sumo $|E_0|$ de ellos. El conjunto de subconjuntos finitos de una base infinita $B$ tiene cardinalidad $|B|$, y una unión de $|B|$ conjuntos finitos tiene cardinalidad a lo sumo $|B|$. Agrupando $S$ por sus soportes obtenemos $|S|\leq|B|$. Usamos aquí el axioma de elección.

Si $S$ también es maximal, intercambiamos sus papeles con $B$ y obtenemos ambas desigualdades; Cantor--Bernstein da la igualdad de cardinales.
\end{proof}

\begin{obs}[Ampliación de escalares y bases infinitas]
El módulo $M$ no es necesariamente un espacio vectorial sobre $K(R)$; por ejemplo, $\mathbb Z$ no lo es sobre $\mathbb Q$. Para ampliar escalares se incluye $M=R^{\oplus E}$ en $K(R)^{\oplus E}$.

En dimensión infinita, una sucesión de intercambios finitos puede perder la generación al pasar al límite. Si $e_0,e_1,\ldots$ es una base, los vectores $s_n=e_n-e_{n+1}$ son independientes. Reemplazar sucesivamente $e_n$ por $s_n$ conserva una base en cada etapa finita, pero los $s_n$ juntos no generan $e_0$: toda combinación finita de ellos tiene suma de coordenadas cero.
\end{obs}

\begin{ejem}
Como espacio vectorial sobre $\mathbb C$, el anillo $\mathbb C[x]$ tiene la base numerable $1,x,x^2,\ldots$. La proposición implica que cualquier subconjunto no numerable de $\mathbb C[x]$ es linealmente dependiente.
\end{ejem}

%%%%%%%%%%%%%%%%%%%%%%%%%%%%%%%%%%%%%%%%%%%%%%%%%%%%%%%%%%%%%%%%%%%%%%%%%%%%%%%
\subsection{Reconstrucción de los índices y definición de rango}
%%%%%%%%%%%%%%%%%%%%%%%%%%%%%%%%%%%%%%%%%%%%%%%%%%%%%%%%%%%%%%%%%%%%%%%%%%%%%%%
% P05.
\begin{ejer}[Remisión a VI.1.8]\label{ejer:C20-indices}
Sea $R$ un dominio. Demostrar que $F^R(A)\cong F^R(B)$ si y sólo si existe una biyección $A\to B$.
\end{ejer}
\begin{sol}
Un isomorfismo transporta bases a bases. La imagen de la base canónica de $F^R(A)$ y la base canónica de $F^R(B)$ tienen igual cardinalidad por la \cref{prop:C20-cardinal}; por tanto $|A|=|B|$.

A la inversa, una biyección $f:A\to B$ induce el homomorfismo $\widetilde f(e_a)=e_{f(a)}$. El inducido por $f^{-1}$ es su inverso, pues las composiciones fijan todos los generadores. El isomorfismo se expresa en el cuadrado:
\[
\begin{tikzcd}[row sep=large,column sep=large]
F^R(A)\arrow[r,"\widetilde f","\sim"']&F^R(B)\\
A\arrow[u,"j_A"]\arrow[r,"f","\sim"']&B\arrow[u,"j_B"'].
\end{tikzcd}
\]
\end{sol}

\begin{coro}
Un dominio satisface la propiedad de número de base invariante: para enteros no negativos $m,n$, se tiene $R^m\cong R^n$ si y sólo si $m=n$.
\end{coro}
\begin{proof}
Aplicamos el ejercicio a conjuntos de $m$ y $n$ elementos.
\end{proof}

\begin{defi}
El \emph{rango} de un módulo libre $M$ sobre un dominio $R$, denotado $\rk_R M$, es la cardinalidad de un subconjunto linealmente independiente maximal. Para un espacio vectorial $V$ sobre $k$ se llama \emph{dimensión} y se escribe $\dim_k V$.
\end{defi}
La proposición garantiza que está bien definido. En un módulo libre, toda base es independiente maximal; sobre un dominio, un independiente maximal puede no ser base, como $\{2\}\subseteq\mathbb Z$, pero tiene la misma cardinalidad que una base. El módulo cero tiene rango cero.

\newpage
%%%%%%%%%%%%%%%%%%%%%%%%%%%%%%%%%%%%%%%%%%%%%%%%%%%%%%%%%%%%%%%%%%%%%%%%%%%%%%%
\section{Torsión y separación de un sumando}
%%%%%%%%%%%%%%%%%%%%%%%%%%%%%%%%%%%%%%%%%%%%%%%%%%%%%%%%%%%%%%%%%%%%%%%%%%%%%%%
%%%%%%%%%%%%%%%%%%%%%%%%%%%%%%%%%%%%%%%%%%%%%%%%%%%%%%%%%%%%%%%%%%%%%%%%%%%%%%%
\subsection{Libertad y ausencia de torsión}
%%%%%%%%%%%%%%%%%%%%%%%%%%%%%%%%%%%%%%%%%%%%%%%%%%%%%%%%%%%%%%%%%%%%%%%%%%%%%%%
% P06--P07.
Desde aquí $R$ es un dominio.
\begin{defi}
Un elemento $m\in M$ es de \emph{torsión} si existe $r\in R\setminus\{0\}$ tal que $rm=0$, es decir, si la familia $\{m\}$ es dependiente. Escribimos
\[
\operatorname{Tor}_R(M)=\set{m\in M\given \exists r\ne0,\ rm=0}.
\]
El módulo es \emph{libre de torsión} si este conjunto es $\{0\}$, y es \emph{de torsión} si coincide con todo $M$.
\end{defi}
Para un anillo conmutativo no nulo, ser libre de torsión como módulo sobre sí mismo equivale a ser dominio: la condición dice precisamente que $rm=0$ y $r\ne0$ implican $m=0$.

\begin{ejer}[Remisión a VI.4.1]\label{ejer:C20-torsion}
Probar que $\operatorname{Tor}_R(M)$ es un submódulo cuando $R$ es un dominio.
\end{ejer}
\begin{sol}
El cero pertenece al conjunto porque $1\ne0$ y $1\cdot0=0$. Si $rm=0$, $sn=0$ con $r,s\ne0$, entonces $rs\ne0$ y
\[
rs(m+n)=s(rm)+r(sn)=0.
\]
Además, para $a\in R$, se tiene $r(am)=a(rm)=0$, lo que incluye a $-m$ tomando $a=-1$. Por tanto el conjunto es cerrado bajo suma y multiplicación por escalares.
\end{sol}

\begin{lema}\label{lema:C20-libre-torsion}
Los submódulos y las sumas directas de módulos libres de torsión son libres de torsión. En particular, todo módulo libre sobre un dominio es libre de torsión.
\end{lema}
\begin{proof}
Una igualdad $rm=0$ en un submódulo vale también en el módulo ambiente, y allí $r\ne0$ fuerza $m=0$. En una suma directa, si $r(m_i)_i=0$ con $r\ne0$, entonces $rm_i=0$ para cada $i$, y todas las coordenadas son cero. El argumento funciona para familias arbitrarias. Finalmente, $R$ es libre de torsión y todo módulo libre es isomorfo a una suma directa de copias de $R$.
\end{proof}

\begin{ejem}\label{ejem:C20-ideal-no-libre}
Un ideal no principal $I$ de un dominio $R$ es libre de torsión pero no es un módulo libre. En efecto, $I\subseteq R$ es libre de torsión. Si tuviera una base $B$, sería independiente en $R$, cuyo rango es uno. Como un ideal no principal es no nulo, su base tendría exactamente un elemento y el ideal sería principal, contradicción.

Consideremos $I=(2,x)\subseteq\mathbb Z[x]$. Para verificar que no es principal, si $I=(f)$ entonces $f\mid2$, luego $f$ es constante; además $f\mid x$ fuerza $f=\pm1$. Pero $I$ es propio: la evaluación en cero seguida de reducción módulo dos anula a $2$ y a $x$, y no anula a $1$. Así, $I$ no puede ser todo el anillo.
\end{ejem}
En particular, un submódulo de un módulo libre sobre un dominio puede no ser libre. Sobre un campo sí lo es, porque todo espacio vectorial admite base.

\begin{obs}
Todo espacio vectorial con base $B$ es isomorfo a $K^{\oplus B}$. Sólo en dimensión finita puede escribirse como $K^n$ con $n$ entero; por ejemplo, $K[x]$ tiene dimensión infinita.
\end{obs}

%%%%%%%%%%%%%%%%%%%%%%%%%%%%%%%%%%%%%%%%%%%%%%%%%%%%%%%%%%%%%%%%%%%%%%%%%%%%%%%
\subsection{Un lema sobre dominios de ideales principales}
%%%%%%%%%%%%%%%%%%%%%%%%%%%%%%%%%%%%%%%%%%%%%%%%%%%%%%%%%%%%%%%%%%%%%%%%%%%%%%%
% P07--P09.
\begin{lema}[Separación de un sumando]\label{lema:C20-separacion}
Sea $R$ un DIP, $F$ un módulo libre finitamente generado y $0\ne M\subseteq F$. Existen $a\ne0$, $x\in F$, $y=ax\in M\setminus\{0\}$ y submódulos $F'\subseteq F$, $M'=F'\cap M$ tales que
\[
F=Rx\oplus F',\qquad M=Ry\oplus M'.
\]
\end{lema}
\begin{proof}
Para $\varphi\in\Hom_R(F,R)$, la imagen $\varphi(M)$ es un ideal de $R$. La colección
\[
\mathcal I=\set{\varphi(M)\given \varphi\in\Hom_R(F,R)}
\]
es no vacía e incluye un ideal no nulo: tomando un elemento no nulo de $M$ y una coordenada no nula en una base de $F$, la proyección correspondiente tiene imagen no nula. Como $R$ es noetheriano, $\mathcal I$ tiene un elemento maximal por inclusión, digamos $\alpha(M)=(a)$, necesariamente no nulo. Elegimos $y\in M$ con $\alpha(y)=a\ne0$.

Veamos que $a\mid\varphi(y)$ para toda $\varphi:F\to R$. Como $R$ es DIP, escribimos
\[
(a,\varphi(y))=(b),\qquad b=ra+s\varphi(y),
\]
y definimos $\psi=r\alpha+s\varphi$. Entonces $\psi(y)=b$, por lo que
\[
\alpha(M)=(a)\subseteq(b)\subseteq\psi(M).
\]
La maximalidad de $\alpha(M)$ en $\mathcal I$ da igualdad de los extremos y, por tanto, $(a)=(b)$. De aquí $\varphi(y)\in(a)$.

Identifiquemos $F$ con $R^n$. Aplicando lo anterior a las proyecciones coordenadas, todas las coordenadas de $y$ son divisibles por $a$. Existe entonces $x\in F$ con $y=ax$. La igualdad
\[
a=\alpha(y)=a\alpha(x)
\]
implica $\alpha(x)=1$, pues $R$ es dominio y $a\ne0$.

Tomemos $F'=\ker\alpha$ y $M'=M\cap F'$. Todo $z\in F$ se escribe
\[
z=\alpha(z)x+\bigl(z-\alpha(z)x\bigr),\qquad
\alpha\bigl(z-\alpha(z)x\bigr)=0.
\]
Además, si $rx\in F'$, entonces $0=\alpha(rx)=r$. Así, $F=Rx\oplus F'$.

Si $z\in M$, tenemos $\alpha(z)=ca$ para algún $c\in R$ y entonces
\[
\alpha(z)x=cy\in Ry,\qquad z-cy\in M\cap\ker\alpha=M'.
\]
Por ello $M=Ry+M'$. Finalmente, si $dy\in M'$, se cumple $0=\alpha(dy)=da$. Como $a\ne0$, resulta $d=0$, y la suma también es directa.
\end{proof}

\begin{obs}[Referencias de lectura]
Para bases, rango y torsión, véase \citet[cap.~VI, 1.1/4--7/9--11/13--14 y 4.1--3]{aluffi2009}. El lema de separación es VI.5.2 y los ejercicios corresponden a VI.1.8 y VI.4.1.
\end{obs}
